\chapter{Centralised Exchanges}\label{ch:m3:centralised-exchanges}

In the week that began on 6 November 2022 the customers of one of the largest crypto exchanges withdrew
seven billion dollars from it. On 8 November it paused all customer withdrawals; on 11 November it filed
for bankruptcy. Its customers had thought of it as a venue: a place where
their orders met other orders. They learned that it had also been their broker, holding their accounts;
their custodian, holding their coins; their clearing house, guaranteeing their trades and liquidating
their leveraged positions; and, through an affiliated trading firm that could borrow from the exchange
without limit, the counterparty that had spent their money. Crypto trading happens mostly on such
venues. This chapter describes what a \omterm{def:m3:centralised-exchanges:cex}{centralised exchange} is, how a trading firm connects to one and
pays for it, and what the firm is exposed to while its assets sit there: the risk that the venue is
not what it says, what proofs of reserves can and cannot establish, and the settlement arrangements
built to keep a firm's collateral out of the venue's hands.

\section{The exchange that is also broker, custodian and clearing house}

In the equity and futures markets of One Quant Book 1, a trade passes through separate institutions,
each regulated for its role: a broker holds the client's account and routes its orders (chapter 4), the
exchange matches them, a central counterparty guarantees the trade and collects margin, and a custodian
or settlement system holds the securities and the cash (chapter 5). Each has its own capital, its own
supervisor and its own rules on what it may do with client assets.

\begin{definition}[Centralised exchange]\label{def:m3:centralised-exchanges:cex}
A \emph{centralised exchange}\index{centralised exchange} is a crypto trading venue run by a company
that holds its customers' assets in wallets and bank accounts it controls, keeps their balances in its
own internal ledger, matches their orders on its own order book, and settles trades by changing that
ledger rather than on a \omterm{def:m3:blockchains-for-traders:chain}{blockchain}.
\end{definition}

Only deposits and withdrawals touch a chain. A customer who sends bitcoin to the exchange gives up the
coins in exchange for a line in the venue's database; every trade after that moves numbers between
lines, and only a withdrawal turns a line back into coins. The same company performs the four roles of
the traditional chain, and usually more: it lends to customers who trade with leverage, runs the
liquidation engine that closes their positions (\cref{ch:m3:margin-liquidation-and-loss-allocation}),
issues its own \omterm{def:m3:blockchains-for-traders:key}{token}, and may operate a trading arm or be affiliated with one.

\begin{omfigure}
\centering
\begin{tikzpicture}[font=\small,
  r/.style={draw,thick,rounded corners=2pt,align=center,minimum height=0.8cm,text width=2.3cm},
  arr/.style={-{Stealth[length=2.2mm]},thick}]
\node[font=\small\bfseries,anchor=west] at (-1.3,2.1) {Securities and futures};
\node[r,draw=omDef,fill=omDef!8] (b) at (0,1.2) {broker\\(accounts)};
\node[r,draw=omThm,fill=omThm!8] (x) at (3.1,1.2) {exchange\\(matching)};
\node[r,draw=omProp,fill=omProp!8] (c) at (6.2,1.2) {central\\counterparty};
\node[r,draw=omExo,fill=omExo!8] (k) at (9.3,1.2) {custodian,\\settlement};
\draw[arr] (b) -- (x); \draw[arr] (x) -- (c); \draw[arr] (c) -- (k);
\node[font=\small\bfseries,anchor=west] at (-1.3,-0.3) {Centralised crypto exchange};
\node[draw=omThm,thick,rounded corners=3pt,fill=omThm!4,minimum width=10.9cm,minimum height=1.2cm,
  align=center] (one) at (4.65,-1.2) {one company: accounts, matching, margin and liquidation,\\
  custody of coins and cash, lending, own token, (often) an affiliated trading arm};
\end{tikzpicture}

\omcaption{In regulated securities and futures markets separate institutions hold the account, match
the order, guarantee the trade and hold the assets; a centralised crypto exchange does all of it inside
one company, whose internal ledger is the only record of who owns what. Schematic.}\label{fig:m3:centralised-exchanges:roles}
\end{omfigure}

\begin{definition}[Customer asset segregation]\label{def:m3:centralised-exchanges:segregation}
\emph{Customer asset segregation}\index{customer asset segregation} is the separation of assets held
for customers from the holder's own assets, in distinct accounts or wallets and under legal terms that
keep them out of the holder's estate if it fails, so that they cannot be used for its own purposes and
return to the customers in an insolvency.
\end{definition}

Segregation is what makes a broker's or a futures commission merchant's failure survivable for its
clients. On a \omterm{def:m3:centralised-exchanges:cex}{centralised exchange} it depends on the venue's terms of service, on the law of the
country where it is incorporated, and on its honesty. Where the terms or the law are silent, a
customer's balance is an unsecured claim on the company, and the coins the venue holds belong to its
creditors as a whole.

For a trading firm, the practical consequence is that a balance on an exchange is a credit exposure to
that exchange, of the full amount, for as long as it sits there. The firm sets limits per venue, as a
bank sets limits per counterparty, and measures what it would lose, and how fast it could withdraw, if
the venue stopped paying.

\section{Accounts, matching and interfaces}

A firm opens a master account with the exchange after an onboarding review (\cref{ch:m3:getting-access-crypto-venues})
and divides it for its strategies and desks.

\begin{definition}[Sub-account, API key]\label{def:m3:centralised-exchanges:subaccount}
A \emph{sub-account}\index{sub-account} is a separately margined account under a master account, with
its own balances, positions and limits, between which the master can move assets. An \emph{API
key}\index{API key} is a credential, a key identifier and a secret, with which a program signs its
requests to an exchange; each key carries permissions (read balances, trade, withdraw) and may be
restricted to a list of IP addresses.
\end{definition}

\omterm{def:m3:centralised-exchanges:subaccount}{Sub-accounts} separate strategies so that a loss or a liquidation in one cannot consume the collateral of
another, and they make each strategy's fees, positions and P\&L visible to the firm without its own
bookkeeping. \omterm{def:m3:centralised-exchanges:subaccount}{API keys} are the firm's most sensitive operational secret after the \omterm{def:m3:blockchains-for-traders:key}{private keys} of its
wallets: a key with withdrawal rights is a key to the money. Trading keys are issued without withdrawal
rights, restricted to the firm's servers' addresses, rotated, and held in a secrets store rather than in
code; withdrawals go through a separate process with its own approvals and address allow-lists.

Matching follows the price-time priority of One Quant Book 1, chapter 19, with differences a trader
notices at once. The market never closes, so there are no opening and closing auctions and no overnight
gap; the order types are fewer, with market, limit and stop orders, their time-in-force flags, and the
\omterm{def:m3:centralised-exchanges:postonly}{post-only order} of the next section; and self-trade prevention, which stops a firm's own buy and sell
orders from matching each other, is a flag on each order with several modes (cancel the resting order,
the incoming one, or both).

A program talks to an exchange over two kinds of interface. Requests (place, amend and cancel an order,
query balances and open orders) go over a request-response interface, usually HTTPS; market data and
the account's own events (fills, balance changes) arrive on streaming WebSocket connections. The
streams' behaviour, their sequence numbers, snapshots and gaps, is the subject of
\cref{ch:m3:crypto-data-and-infrastructure}. The request side is where the venue protects itself.

\begin{definition}[Rate limit, request weight]\label{def:m3:centralised-exchanges:ratelimit}
A \emph{rate limit}\index{rate limit} is the maximum number of requests, or of units of load, that a
venue accepts from a client in a time interval, beyond which it rejects requests and may ban the client
for a period. A \emph{request weight}\index{request weight} is the number of units a request consumes
against such a limit, set by the venue per endpoint to reflect its cost: a full order-book snapshot
weighs more than a single order.
\end{definition}

\begin{dated}{2026-09}{Rate limits at a large venue}\label{dat:m3:centralised-exchanges:limits}
Binance's spot API documentation assigns each endpoint a weight and returns the weight used so far in
each interval in a response header; limits apply per IP address, not per \omterm{def:m3:centralised-exchanges:subaccount}{API key}, and order-count
limits per account, also reported in a header. A request that breaks a limit receives HTTP status 429
with a \texttt{Retry-After} header; an IP that keeps sending after 429 responses is automatically
banned, with status 418, for periods that scale for repeat offenders from 2 minutes to 3 days. On 24
September 2026 its exchange-information endpoint returned limits of 6\,000 units of \omterm{def:m3:centralised-exchanges:ratelimit}{request weight} per
minute, 100 orders per 10 seconds, 200\,000 orders per day and 300\,000 raw requests per 5 minutes;
they change, and a client reads them rather than hard-coding them.
\end{dated}

Two properties of such limits shape a client's design. The first is that the windows are fixed.

\begin{proposition}[The burst at a window boundary]\label{prop:m3:centralised-exchanges:burst}
Under a limit of $L$ units per fixed window of length $T$ aligned on the clock, a client can send $2L$
units within an interval of length shorter than $T$ that straddles a window boundary, but never more
than $2L$ in any interval of length $T$. A sliding window, which counts the units of the last $T$
seconds, caps every interval of length $T$ at $L$.
\end{proposition}

\begin{proof}
Send $L$ units just before the boundary and $L$ just after: each window holds $L$. Any interval of
length $T$ meets at most two windows, each holding at most $L$. A sliding window counts, at each
instant, exactly the units sent in the preceding interval of length $T$ and refuses any unit that would
take it above $L$.
\end{proof}

A client that relies on the boundary burst gets an advantage for a moment and is sometimes caught by a
venue that counts differently from what it documents. The robust design counts conservatively on its own
side, never sends a request it expects to be refused, and treats a 429 as a bug to be fixed, since the
ban that follows repeated ones takes the firm off the venue.

The second property matters more for a market maker. Every quote change is an order, or a cancel and an
order, and order counts are limited per account. When the market moves fast the strategy wants to change
its quotes most often, which is exactly when the limit binds. A client that queues the updates it may not
yet send sends them later, in order: quotes computed for a market that has moved.
\Cref{fig:m3:centralised-exchanges:burst} shows the effect.

\begin{omfigure}
\begin{tikzpicture}
\begin{axis}[width=11cm,height=5cm,axis y line*=left,
  xlabel={second},ylabel={queued quote updates},
  tick label style={font=\small},label style={font=\small},
  xmin=0,xmax=299,ymin=0,ymax=1000,grid=major,grid style={gray!25},
  table/col sep=comma]
\addplot[omThm,very thick] table[x=second,y=fifo]{figdata/markets-3/15-centralised-exchanges/burst.csv};\label{plot:m3:burstfifo}
\addplot[omExo,very thick,dashed] table[x=second,y=coalesced]{figdata/markets-3/15-centralised-exchanges/burst.csv};\label{plot:m3:burstcoal}
\end{axis}
\begin{axis}[width=11cm,height=5cm,axis y line*=right,axis x line=none,
  ylabel={oldest queued update, s},ylabel near ticks,
  tick label style={font=\small},label style={font=\small},
  xmin=0,xmax=299,ymin=0,ymax=100,
  legend columns=3,legend style={font=\footnotesize,at={(0.5,-0.3)},anchor=north,draw=none,fill=none,/tikz/every even column/.append style={column sep=6pt}},
  table/col sep=comma]
\addlegendimage{/pgfplots/refstyle=plot:m3:burstfifo}\addlegendentry{queued, sent in order (left)}
\addlegendimage{/pgfplots/refstyle=plot:m3:burstcoal}\addlegendentry{coalesced (left)}
\addplot[omDef,thick] table[x=second,y=age]{figdata/markets-3/15-centralised-exchanges/burst.csv};\addlegendentry{age of oldest (right)}
\end{axis}
\end{tikzpicture}

\omcaption{A quoting strategy on five instruments under a limit of 100 orders per 10 seconds: 5 quote
updates a second, 25 a second in a volatile minute (seconds 60 to 119). Queued and sent in order, the
backlog reaches 900 updates and the oldest waits 89 seconds; coalesced (only the latest quote per
instrument kept), nothing waits. Illustrative rates; the limits are those of \cref{dat:m3:centralised-exchanges:limits}.
Data: the chapter's tutorial.}\label{fig:m3:centralised-exchanges:burst}
\end{omfigure}

The cure is to coalesce: keep, per instrument and side, only the latest desired quote, and send the
difference between it and what is resting when the limit allows. The backlog is then at most one update
per instrument and side, and what is sent is current. The build of this chapter is the governor that
decides, before each request, whether the venue will accept it.

\section{Fees, rebates and market-maker programmes}

Crypto venues charge by maker-taker pricing with volume tiers (One Quant Book 1, chapter 4): a taker,
whose order executes on arrival, pays more than a maker, whose resting order provides the liquidity, and
both pay less as their thirty-day volume grows. A strategy that intends to provide liquidity must be sure
its orders never take it by accident.

\begin{definition}[Post-only order]\label{def:m3:centralised-exchanges:postonly}
A \emph{post-only order}\index{post-only order} is a limit order that the venue accepts only if it would
rest on the book; if it would execute at once against a resting order, the venue rejects it (or, on
some venues, reprices it one tick away) instead of filling it as a taker.
\end{definition}

A \omterm{def:m3:centralised-exchanges:postonly}{post-only order} guarantees the maker fee and the absence of a crossing trade, at the price of a
rejection when the market has moved through the order's price while it travelled. For a market maker
whose quotes are computed from a price a few milliseconds old, that rejection is the right outcome.

\begin{dated}{2026-09}{Spot fee schedules at two venues}\label{dat:m3:centralised-exchanges:fees}
Binance's spot fee schedule charges a regular user 0.100\% as maker and as taker (0.075\% when paid in
its own \omterm{def:m3:blockchains-for-traders:key}{token}, BNB), and its highest tier, VIP~9 (thirty-day volume of at least USD~4 billion or a BNB
holding), 0.011\% as maker and 0.023\% as taker. Kraken Pro charges its lowest spot tier 0.40\% as maker
and 0.80\% as taker, and its highest tier 0.00\% and 0.05\%. Schedules change often, and large market
makers negotiate terms outside them.
\end{dated}

\begin{table}[ht]
\centering
\small
\begin{tabular}{@{}lrrrr@{}}
\toprule
Venue and tier & maker (bp) & taker (bp) & both legs taker (bp) & both legs maker (bp) \\
\midrule
Binance, regular & 10.0 & 10.0 & 20.0 & 20.0 \\
Binance, VIP 9 & 1.1 & 2.3 & 4.6 & 2.2 \\
Kraken Pro, lowest & 40.0 & 80.0 & 160.0 & 80.0 \\
Kraken Pro, highest & 0.0 & 5.0 & 10.0 & 0.0 \\
\bottomrule
\end{tabular}
\caption{The cost of a round trip (a buy and a sale) in basis points of notional, from the schedules of
\cref{dat:m3:centralised-exchanges:fees}. A strategy whose expected edge per round trip is a few basis
points exists only at the top tiers.}\label{tab:m3:centralised-exchanges:fees}
\end{table}

\Cref{tab:m3:centralised-exchanges:fees} makes the point of the tiers. A strategy that captures five
basis points per round trip loses money at the lowest tier of either venue and makes it at the highest;
the difference is not skill but volume. The venues' market-maker programmes go further, paying rebates
(negative maker fees) and sometimes a fixed fee in return for quoting obligations: a maximum spread, a
minimum size, a share of the time. These are incentive programmes in the sense of One Quant Book 1, and
\cref{ch:m3:getting-access-crypto-venues} works through one. Two cautions apply. A discount paid in the
venue's own \omterm{def:m3:blockchains-for-traders:key}{token} is a position in that \omterm{def:m3:blockchains-for-traders:key}{token}, whose value depends on the venue's health. And a
volume-based schedule rewards volume, including volume that trades with itself, which is one reason why
reported volumes on some venues overstate real activity (\cref{ch:m3:spot-markets}).

\section{The exchange as counterparty: the failure of FTX}

The customers of FTX discovered the difference between a venue and a counterparty in a week. The account
below comes from the charges brought by the US Securities and Exchange Commission and from the filings of
the debtors, the administrators who took over the companies in bankruptcy.

FTX ran a large international exchange, FTX.com, and a smaller US one, FTX.US. Alameda Research, a
trading firm owned by the same founder, came to dominate borrowing on FTX.com. In December 2022
the SEC charged the founder, Samuel Bankman-Fried, with defrauding the equity investors from whom FTX had
raised more than USD~1.8 billion, by concealing that FTX customers' funds were being diverted to Alameda,
and that the commingled customer money was used for venture investments, real estate and political
donations. The debtors' reports describe how. The exchange's code gave Alameda a borrowing limit of USD~65
billion and settings that let it withdraw assets and hold a negative balance without being liquidated,
privileges no other customer had. Customer deposits of dollars were partly received into accounts
controlled by Alameda. Almost all crypto assets were held in internet-connected wallets, without the
multi-signature controls common in the industry.

\begin{omfigure}
\begin{tikzpicture}
\begin{axis}[width=11cm,height=5cm,ybar,bar width=9pt,
  xlabel={November 2022},ylabel={USD million},
  tick label style={font=\small},label style={font=\small},
  xmin=0.4,xmax=11.6,xtick={1,2,3,4,5,6,7,8,9,10,11},ymin=-500,ymax=8000,grid=major,grid style={gray!25},
  legend columns=3,legend style={font=\footnotesize,at={(0.5,-0.3)},anchor=north,draw=none,fill=none,/tikz/every even column/.append style={column sep=10pt}},
  table/col sep=comma]
\addplot[fill=omThm!60,draw=omThm,area legend] table[x=day,y=cust_out]{figdata/markets-3/15-centralised-exchanges/flows.csv};\addlegendentry{customers' net withdrawals}
\addplot[omDef,very thick,sharp plot,mark=*,mark size=1.5pt] table[x=day,y=cust_cum]{figdata/markets-3/15-centralised-exchanges/flows.csv};\addlegendentry{cumulative}
\addplot[omExo,very thick,sharp plot,dashed,mark=square*,mark size=1.5pt] table[x=day,y=rel_cum]{figdata/markets-3/15-centralised-exchanges/flows.csv};\addlegendentry{related parties' net deposits, cum.}
\end{axis}
\end{tikzpicture}

\omcaption{FTX.com, 1 to 11 November 2022: customers withdrew, net, USD~7.0 billion, 3.3 billion of it on
7 November; Alameda and other related parties deposited, net, USD~3.2 billion, mostly on 6 to 8 November.
USD millions at petition-time prices, as published; cumulative sums differ from the filing's totals by
rounding. Data: FTX Debtors, presentation of 2 March 2023, In re FTX Trading Ltd., Bankr.\ D.\ Del.\ No.\
22-11068.}\label{fig:m3:centralised-exchanges:flows}
\end{omfigure}

When doubts about Alameda's balance sheet became public in early November 2022, customers ran.
\Cref{fig:m3:centralised-exchanges:flows} shows the run in the debtors' figures: net customer withdrawals
from FTX.com of about USD~1.8 billion on 6 November and 3.3 billion on 7 November, before withdrawals
effectively stopped. The companies filed for Chapter~11 bankruptcy on 11 November (and some affiliates on
14 November). On the day of the filing, unauthorised transfers drained a further few hundred million
dollars from the exchanges' wallets.

\begin{table}[ht]
\centering
\small
\begin{tabular}{@{}lrrr@{}}
\toprule
FTX.com, USD million & customer claims & assets located & coverage \\
\midrule
Cash and \omterm{def:m3:blockchains-for-traders:stable}{stablecoins} & 6\,991 & 270 & 3.9\% \\
Bitcoin & 1\,591 & 1 & 0.1\% \\
Ether & 922 & 9 & 1.0\% \\
All liquid \omterm{def:m3:blockchains-for-traders:key}{tokens} (Category A) & 10\,544 & 694 & 6.6\% \\
All \omterm{def:m3:blockchains-for-traders:key}{tokens}, with receivables & 11\,233 & 2\,540 & 22.6\% \\
\bottomrule
\end{tabular}
\caption{Customer claims on FTX.com and the assets the debtors located in its wallets at the petition time,
USD millions at petition-time prices. Category A: \omterm{def:m3:blockchains-for-traders:key}{tokens} with a market capitalisation of at least USD~15
million and a daily volume of at least USD~1 million; most Category B assets were \omterm{def:m3:blockchains-for-traders:key}{tokens} linked to FTX and
Alameda. Receivables are what customers owed the exchange. Preliminary figures, from the debtors'
presentation of 2 March 2023.}\label{tab:m3:centralised-exchanges:shortfall}
\end{table}

\Cref{tab:m3:centralised-exchanges:shortfall} is the balance sheet customers were in fact exposed to. Of
USD~10.5 billion of claims in liquid assets, the exchange held 694 million. It held one million dollars of
bitcoin against 1.6 billion owed. The only assets worth more than the claims on them were illiquid \omterm{def:m3:blockchains-for-traders:key}{tokens}
related to the group itself, valued at prices no one could sell at. The debtors' later complaint against
the founders records that by August 2022 they privately estimated the exchange owed customers more than
USD~8 billion in dollars it could not repay.

\begin{dated}{2026-09}{FTX afterwards}\label{dat:m3:centralised-exchanges:ftx}
Samuel Bankman-Fried was convicted of fraud by a jury in New York in November 2023 and sentenced on 28
March 2024 to 25 years in prison, with a forfeiture of more than USD~11 billion. On 8 August 2024 a
federal court, in a case brought by the Commodity Futures Trading Commission, ordered FTX and Alameda to
pay USD~8.7 billion in restitution and 4 billion in disgorgement, finding that FTX had claimed to
segregate customer assets while commingling and misappropriating them. Under the US Bankruptcy Code a claim is
determined in dollars as of the date of the petition, so a customer owed a bitcoin was owed its November
2022 dollar value: the plan confirmed on 7 October 2024 promised 98\% of creditors by number about 119\% of
their allowed claims, in dollars, whatever the coins had since become worth.
\end{dated}

FTX was not the first such failure. Mt.~Gox, an early bitcoin exchange, filed for bankruptcy in Tokyo
on 28 February 2014, saying it might have lost about 750\,000 of its customers' bitcoins and 100\,000 of
its own; the court converted them into a civil rehabilitation in June 2018, and the
trustee began repaying creditors in bitcoin and bitcoin cash in 2024, ten years after the failure. The
lessons for a trading firm are practical. A balance on a venue is an unsecured loan to it unless the firm
knows otherwise; the loan should be sized as one, spread across venues, and swept back to the firm's own
custody when it is not needed for trading. A run starts before the news is certain, and the first to
withdraw are paid. And a venue's own statements about its reserves are a claim, which the next section
examines.

\section{Proof of reserves and its limits}

After November 2022 several exchanges began to publish evidence that they held their customers' assets.

\begin{definition}[Proof of reserves, proof of liabilities]\label{def:m3:centralised-exchanges:por}
A \emph{proof of reserves}\index{proof of reserves} is evidence that an exchange controls on-chain assets
at least equal to its customers' balances at a moment: signatures from its addresses, whose balances
anyone can read on the chains, compared with a total of customer balances. A \emph{proof of
liabilities}\index{proof of liabilities} is the evidence that this total is complete and correct: a
commitment to every customer balance, usually the root of a Merkle tree (a binary tree in which each node
is the hash of its children), against which each customer can check that her own balance was counted.
\end{definition}

In a Merkle sum tree each leaf commits to one account, by hashing the account's identifier, a secret nonce
known to its owner, and its balance, and carries the balance; each parent hashes its two children together
with the sum of their balances and carries that sum. The root's hash commits to every leaf and its sum is
the published total of liabilities. A customer receives her leaf's nonce and the siblings on the path from
her leaf to the root, one per level, and recomputes the root: $\log_2 n$ hashes for $n$ accounts, eleven
for two thousand.

\begin{omfigure}
\centering
\begin{tikzpicture}[font=\footnotesize,
  n/.style={draw,thick,rounded corners=2pt,align=center,minimum height=0.75cm,text width=1.35cm},
  e/.style={thick,gray}]
\node[n,draw=omThm,fill=omThm!12] (r) at (5.25,3.3) {root\\$\Sigma=90$};
\node[n,draw=omExo,fill=omExo!10] (a) at (2.6,2.1) {$h_{12}$\\$\Sigma=35$};
\node[n,draw=omThm,fill=omThm!12] (b) at (7.9,2.1) {$h_{34}$\\$\Sigma=55$};
\node[n] (l1) at (1.3,0.8) {acct 1\\30};
\node[n] (l2) at (3.9,0.8) {acct 2\\5};
\node[n,draw=omExo,fill=omExo!10] (l3) at (6.6,0.8) {acct 3\\40};
\node[n,draw=omThm,fill=omThm!12] (l4) at (9.2,0.8) {acct 4\\15};
\draw[e] (r) -- (a); \draw[e] (r) -- (b); \draw[e] (a) -- (l1); \draw[e] (a) -- (l2);
\draw[e] (b) -- (l3); \draw[e] (b) -- (l4);
\node[anchor=north,align=center,text width=11cm] at (5.25,0.2) {Account 4 recomputes its leaf, then
$h_{34}$ from the sibling leaf of account 3, then the root from the sibling node $h_{12}$ (shaded
grey); it checks every sum on the way is non-negative.};
\end{tikzpicture}

\omcaption{A Merkle sum tree of four accounts. Each node carries the hash of its children and the sum of
their balances; the root commits to all balances and their total. A customer's inclusion proof is the
siblings on her path (grey for account 4). Schematic.}\label{fig:m3:centralised-exchanges:merkle}
\end{omfigure}

\begin{proposition}[What an inclusion proof establishes]\label{prop:m3:centralised-exchanges:inclusion}
Let the hash function be collision-resistant. (i) If a customer's leaf and path recompute the published
root, her balance is a leaf of the committed tree and is counted in its total. (ii) If moreover every leaf
balance is non-negative, the published total is at least the sum of the balances of all customers whose
proofs verify. Without (ii), a tree containing a negative leaf verifies for every customer while its total
understates the liabilities by the size of that leaf.
\end{proposition}

\begin{proof}
(i) A different leaf or path reaching the same root would give two different inputs with the same hash at
some level. (ii) The root's sum is the sum of all leaves; removing the non-negative leaves of the customers
who did not check can only lower it. For the last claim, a negative leaf changes no other leaf's hash, so
every honest path still recomputes the root, whose sum is lowered by the negative amount.
\end{proof}

The negative leaf is not a curiosity: it is the cheapest way for an insolvent exchange to shrink its
published liabilities to match its reserves. It is caught when customers' verifiers check that every sum
on their path is non-negative, since the negative leaf makes its ancestors' sums negative until a subtree
holds more real balances than it hides. The tutorial builds a tree of 1\,023 customers and one fake
account that hides 60\% of the liabilities: all 1\,023 proofs pass a hash-only check, and all 1\,023 fail
a check of the sums. Schemes that do not want to reveal sibling sums use zero-knowledge range proofs to show
that every leaf is non-negative. Omitting accounts is harder to catch: it is found only when an omitted
customer checks, so an exchange that omits the dormant accounts least likely to check is caught with small
probability.

\begin{dated}{2026-09}{A published proof of reserves}\label{dat:m3:centralised-exchanges:kraken}
Kraken describes its \omterm{def:m3:centralised-exchanges:por}{proof of reserves} as an independent accountant's review: an anonymised snapshot of
client balances aggregated in a Merkle tree, each leaf the first 16 characters of a SHA-256 hash of a record
identifier and the balances, and a comparison of those balances with on-chain holdings whose control the
accountant verifies by signatures. Clients can verify their inclusion in the tree themselves. Kraken lists
what the review does not prove: exclusive possession of the \omterm{def:m3:blockchains-for-traders:key}{private keys}, the absence of hidden liabilities
or of funds borrowed to pass the review, off-chain assets, and anything after the snapshot; the assets in
scope were BTC, ETH, SOL, USDC, USDT and XRP.
\end{dated}

A \omterm{def:m3:centralised-exchanges:por}{proof of reserves} is thus a snapshot of one side of a balance sheet with a partial check of the other.
It cannot show liabilities the exchange owes to lenders rather than to customers, assets lent out the day
after, or keys held by someone else too. It is useful evidence against the crudest fraud, and it would not
have revealed the FTX borrowing privilege, which lived in the matching engine's settings, not in the
balances.

\section{Off-exchange settlement}

If the risk is that the venue holds the collateral, the remedy is for it not to.

\begin{definition}[Off-exchange settlement]\label{def:m3:centralised-exchanges:oes}
\emph{Off-exchange settlement}\index{off-exchange settlement} is an arrangement in which a trading firm's
collateral stays with an independent custodian, the exchange credits the firm with a trading balance
mirroring it, and the gains and losses of the firm's trading are settled between the custodian and the
exchange at set intervals or on demand, so that only the unsettled amount is at risk to the exchange.
\end{definition}

The arrangement moves the firm's exposure from the exchange's balance sheet to the custodian's, which is
chosen for its segregation, controls and supervision, and limits what an exchange failure can take to one
settlement period's net gains. It also lets a firm trade on several venues against one pool of collateral,
netting across them. It is not free of risk: the custodian is now the concentration, the exchange extends
credit to the firm between settlements and prices that credit into its limits, and the firm needs its
positions, its mirrored balances and the settlements to reconcile every cycle.

\begin{dated}{2026-09}{An off-exchange settlement network}\label{dat:m3:centralised-exchanges:clearloop}
Copper, a custodian, describes its ClearLoop network as ``\omterm{def:m3:centralised-exchanges:oes}{off-exchange settlement}'' with ``single net
settlement'' across connected venues, which it lists as including OKX, Coinbase, Bybit, Gate.io, Deribit,
Bitget, Bitfinex and Kraken; assets are held in Copper's custody. Other custodians and some exchanges offer
similar arrangements; terms, settlement frequency and connected venues are the providers' and change.
\end{dated}

\begin{omfigure}
\centering
\begin{tikzpicture}[font=\small,
  p/.style={draw,thick,rounded corners=2pt,align=center,minimum height=1cm,text width=2.7cm},
  arr/.style={-{Stealth[length=2.2mm]},thick}]
\node[p,draw=omDef,fill=omDef!8] (f) at (0,0) {trading firm};
\node[p,draw=omProp,fill=omProp!8] (c) at (5,0) {independent\\custodian\\(collateral)};
\node[p,draw=omThm,fill=omThm!8] (x1) at (10,0.9) {exchange A\\(mirrored balance)};
\node[p,draw=omThm,fill=omThm!8] (x2) at (10,-0.9) {exchange B\\(mirrored balance)};
\draw[arr] (f) -- node[above,font=\footnotesize] {deposits} (c);
\draw[arr] (c) -- node[above,sloped,font=\footnotesize] {credit line} (x1);
\draw[arr] (c) -- node[below,sloped,font=\footnotesize] {credit line} (x2);
\draw[arr,dashed] (f.north) to[bend left=25] node[above,font=\footnotesize] {orders} (x1.north);
\draw[arr,dashed] (f.south) to[bend right=25] node[below,font=\footnotesize] {orders} (x2.south);
\end{tikzpicture}

\omcaption{\omterm{def:m3:centralised-exchanges:oes}{Off-exchange settlement}: the firm's collateral stays with a custodian, the exchanges credit a
mirrored trading balance, and net gains and losses are settled between custodian and exchanges at set times.
An exchange's failure can take only the unsettled amount. Schematic.}\label{fig:m3:centralised-exchanges:oes}
\end{omfigure}

\section{Tutorial: checking a proof of reserves}\label{tut:m3:centralised-exchanges:por}

\begin{tutorial}
\textbf{Goal.} Build a Merkle sum tree of customer balances, produce and verify inclusion proofs, and
construct the counter-example of \cref{prop:m3:centralised-exchanges:inclusion}. \textbf{End state:} the
numbers of the section on \omterm{def:m3:centralised-exchanges:por}{proof of reserves}.
\begin{tutsteps}
\item \textbf{The tree.} Leaves commit to account, nonce and balance; parents to their children and the sum.
\omcode{code/markets-3/15-centralised-exchanges/python/m3_cex.py}{14}{52}{A Merkle sum tree, inclusion proofs and their verification.}\label{lst:m3:centralised-exchanges:tree}
\item \textbf{The counter-example.} Insert one account with a negative balance equal to 60\% of the
liabilities and count the customers whose proofs pass with and without the check of sums:
\texttt{fake\_negative(customers())}.
\item \textbf{The shortfall.} \texttt{ftx\_coverage()} reads the debtors' table and returns the shares of
\cref{tab:m3:centralised-exchanges:shortfall}; \texttt{fig\_cex.py} writes the chapter's chart data.
\end{tutsteps}
\textbf{What to change next.} Hide only 2\% of the liabilities and count the customers who still detect it;
put the fake leaf next to the largest account; replace the sibling sums by commitments and ask what a
customer can still check.
\end{tutorial}

\section{Build: the rate-limit governor}\label{bld:m3:centralised-exchanges:ratelimit}

\begin{build}
\bfield{Purpose} Every request the miniature firm sends to a venue passes a governor that knows the venue's
limits and refuses, before sending, any request the venue would reject; after a 429 or a 418 it stops for
the time the venue says. A trading system that learns about limits from the venue's rejections gets banned.
\bfield{Interface} \texttt{Rule(kind, interval\_ms, limit)} for weight or order-count rules;
\texttt{Governor(rules)}; \texttt{try\_send(now, weight, is\_order)} returns whether to send and, if not, how
long to wait; \texttt{on\_status(now, status, retry\_after\_ms)}. The C++20 header \texttt{firm\_ratelimit.hpp}
and the Rust crate \texttt{firm\_ratelimit} implement the same interface; the Python module is the reference.
\bfield{Rules} Fixed windows aligned on the epoch, as the venue counts; a request is sent only if every rule
has room, and then consumes all of them; a blocked governor refuses everything until the block ends; integer
milliseconds, no floating point.
\bfield{Acceptance tests} \texttt{code/firm/ratelimit/tests/}, \texttt{cpp/firm\_ratelimit\_test.cpp} and the
Rust crate's tests: a refused request's waiting time to the next window; order and weight rules counted
separately; backoff after 429 and 418.
\bfield{Stretch} Sliding windows; per-endpoint weights read from the venue's exchange-information response;
reconciliation with the venue's usage headers; a coalescing quote sender on top of the governor.
\end{build}

\omcode{code/firm/ratelimit/cpp/firm_ratelimit.hpp}{41}{60}{The C++20 governor's decision and its reaction to the venue's status codes.}\label{lst:m3:centralised-exchanges:governor}

\begin{omsources}
\item US Securities and Exchange Commission, ``SEC Charges Samuel Bankman-Fried with Defrauding Investors in
Crypto Asset Trading Platform FTX'', press release 2022-219, 13 December 2022.
\item FTX Debtors, ``Preliminary Analysis of Shortfalls'', presentation to the Official Committee of Unsecured
Creditors, In re FTX Trading Ltd., No.\ 22-11068 (JTD), Bankr.\ D.\ Del., Doc 792-1, 2 March 2023.
\item J.~J.~Ray III, \emph{First Interim Report to the Independent Directors on Control Failures at the FTX
Exchanges}, Doc 1242-1, 9 April 2023; FTX Debtors, complaint against the founders, Doc 1886, 20 July 2023.
\item Securities and Exchange Commission v.\ Bankman-Fried, complaint, S.D.N.Y., 13 December 2022; FTX Debtors, press
release on the confirmation of the plan, 7 October 2024; NPR, report of Mt.~Gox's filing, 28 February 2014.
\item Commodity Futures Trading Commission, press release 8938-24, 8 August 2024; UPI, report of the
sentencing, 28 March 2024.
\item Binance, Spot API documentation (limits, enums) and fee schedule; Kraken, fee schedule and proof of
reserves pages; Copper, ClearLoop product page; MtGox rehabilitation trustee's website. All accessed
September 2026.
\end{omsources}

\section{Exercises}

\begin{exercise}[$\star$]\label{exo:m3:centralised-exchanges:1}
Name the four roles that separate institutions play in a regulated futures trade and that a \omterm{def:m3:centralised-exchanges:cex}{centralised
exchange} plays alone.
\end{exercise}

\begin{exercise}[$\star$]\label{exo:m3:centralised-exchanges:2}
A limit is 1\,200 units of weight per minute in fixed windows. What is the most a client can send in any
two seconds? In any minute?
\end{exercise}

\begin{exercise}[$\star$]\label{exo:m3:centralised-exchanges:3}
Why should an \omterm{def:m3:centralised-exchanges:subaccount}{API key} used by a trading strategy not carry withdrawal rights?
\end{exercise}

\begin{exercise}[$\star\star$]\label{exo:m3:centralised-exchanges:4}
A strategy trades USD~50 million a day, half as maker and half as taker, at Binance's regular tier and then
at VIP~9 (\cref{dat:m3:centralised-exchanges:fees}). What does it pay in fees a day in each case?
\end{exercise}

\begin{exercise}[$\star\star$]\label{exo:m3:centralised-exchanges:5}
An exchange has 2\,000\,000 accounts. How many hashes does a customer compute to verify her inclusion? If
the exchange omits 100 accounts and each customer checks with probability 1\%, what is the probability that
at least one omission is found?
\end{exercise}

\begin{exercise}[$\star\star$]\label{exo:m3:centralised-exchanges:6}
From \cref{tab:m3:centralised-exchanges:shortfall}, what share of cash and \omterm{def:m3:blockchains-for-traders:stable}{stablecoin} claims could the
exchange have paid from the cash, \omterm{def:m3:blockchains-for-traders:stable}{stablecoins} and receivables of that line (USD~270 and 310 million)?
\end{exercise}

\begin{exercise}[$\star\star\star$]\label{exo:m3:centralised-exchanges:7}
\emph{Coding.} With \texttt{fake\_negative}, hide 2\% of the liabilities instead of 60\%. How many of the
1\,023 customers detect it when they check sums? Explain why fewer do.
\end{exercise}

\begin{exercise}[$\star\star\star$]\label{exo:m3:centralised-exchanges:8}
\emph{Find the flaw.} ``The exchange published a \omterm{def:m3:centralised-exchanges:por}{proof of reserves} last month showing assets above customer
balances, so our USD~40 million there is safe.''
\end{exercise}

\section{Problem: The Hole in the Balance Sheet}

\begin{problem}[{Weekend problem --- what a customer of FTX.com was owed and what was there}]\label{pb:m3:centralised-exchanges:1}
Use \cref{tab:m3:centralised-exchanges:shortfall}, \cref{fig:m3:centralised-exchanges:flows} and the
debtors' figures quoted in the text. All amounts in USD millions at petition-time prices.

\medskip\noindent\textbf{Part I --- The run.}
\begin{enumerate}
\item How much did customers withdraw, net, from FTX.com from 1 to 11 November 2022?
\item On which day was the largest net withdrawal, and how large was it?
\item What did related parties do over the same days, and why might that matter to customers?
\item Why did withdrawals almost stop from 9 November?
\item Who was paid in full in this run, and why?
\end{enumerate}

\medskip\noindent\textbf{Part II --- The shortfall.}
\begin{enumerate}[resume]
\item What share of claims in liquid \omterm{def:m3:blockchains-for-traders:key}{tokens} was covered by the liquid \omterm{def:m3:blockchains-for-traders:key}{tokens} located?
\item What share with the receivables of those \omterm{def:m3:blockchains-for-traders:key}{tokens} (USD~383 million)?
\item What share of bitcoin claims was covered?
\item What share of all claims was covered by all assets and receivables, and what was the deficit?
\item Why is the last share misleading?
\end{enumerate}

\medskip\noindent\textbf{Part III --- What would have shown it.}
\begin{enumerate}[resume]
\item Would a Merkle-tree \omterm{def:m3:centralised-exchanges:por}{proof of liabilities} have revealed the shortfall?
\item Would a \omterm{def:m3:centralised-exchanges:por}{proof of reserves}, with on-chain holdings compared against the total, have revealed it?
\item Would either have revealed Alameda's borrowing privilege?
\item How would \omterm{def:m3:centralised-exchanges:oes}{off-exchange settlement} have changed a firm's loss?
\item What should a firm's venue limit have been, measured how?
\end{enumerate}

\medskip\noindent\textbf{Part IV --- Judgement.}
\begin{enumerate}[resume]
\item Why were claims fixed in dollars at petition-time prices, and who gained or lost by it?
\item How is a customer balance on a \omterm{def:m3:centralised-exchanges:cex}{centralised exchange} different from an account at a broker with
segregated client assets?
\item What signs, before 6 November, might a risk manager have acted on?
\item State the \emph{named result}: the share of FTX.com customer claims in liquid \omterm{def:m3:blockchains-for-traders:key}{tokens} covered by liquid
\omterm{def:m3:blockchains-for-traders:key}{tokens} located at the petition time, without and with receivables.
\item In one sentence: what is a balance on a \omterm{def:m3:centralised-exchanges:cex}{centralised exchange}?
\end{enumerate}
\end{problem}

\section{Interview questions}

\begin{interviewq}[$\star$ \iqroles{trader}]\label{iq:m3:centralised-exchanges:1}
What is a \omterm{def:m3:centralised-exchanges:postonly}{post-only order} and when would you use one?
\end{interviewq}

\begin{interviewq}[$\star$ \iqroles{risk}]\label{iq:m3:centralised-exchanges:2}
How would you limit a trading firm's exposure to the failure of a crypto exchange?
\end{interviewq}

\begin{interviewq}[$\star\star$ \iqroles{developer}]\label{iq:m3:centralised-exchanges:3}
Your system receives HTTP 429 from an exchange several times an hour. What do you do?
\end{interviewq}

\begin{interviewq}[$\star\star$ \iqroles{developer, researcher}]\label{iq:m3:centralised-exchanges:4}
Explain how a customer checks her inclusion in an exchange's \omterm{def:m3:centralised-exchanges:por}{proof of liabilities}, and how an exchange could
cheat it.
\end{interviewq}

\begin{interviewq}[$\star\star$ \iqroles{trader, risk}]\label{iq:m3:centralised-exchanges:5}
What does a \omterm{def:m3:centralised-exchanges:por}{proof of reserves} not prove?
\end{interviewq}

\begin{interviewq}[$\star\star\star$ \iqroles{developer}]\label{iq:m3:centralised-exchanges:6}
Design the component that sends quote updates for 200 instruments to a venue that allows 100 orders per 10
seconds per account.
\end{interviewq}
